% 引言：研究问题、四项研究与章节的对应、问题速答、记号、证据标签、方法与局限
\chapter*{引言}
\addcontentsline{toc}{chapter}{引言}
\markboth{引言}{引言}

几乎每一款有排位的对战游戏里，都流传着“ELO 机制”的说法：连赢几把就会被塞进“坑队友”的队伍，系统要把每个人的胜率拉回五成，低分段的人“出不去”。这些说法把两件不同的事混在一起：系统怎样估计一个人有多强（评分），以及系统怎样安排谁和谁打（匹配）。本文把这个话题拆成四项研究，分别回答“有什么研究和披露”“胜率能不能算”“诚实的系统会不会让人有被操纵的感觉”“从训练模型的角度看，目标究竟是胜率还是实力”。

\section*{研究问题}

本研究的问题由一位在做酒馆战棋 AI 的提问者（下称“提问者”）提出，原始问题如下（原文照录）：

\begin{quote}\small
我想看看一项研究的情况：各种游戏中的elo机制。请分4个研究进行。\\
1.首先是文献调研，通用的elo或者匹配机制有没有人研究？各个游戏专门的呢？游戏公司有披露的相关机制吗？有公司的专利或者论文吗？\\
2.然后是对于一个游戏，假设结果只有输或者赢，在不干涉任何匹配的情况下，能数学计算出胜率吗？和什么有关呢？比如一把游戏的人数，匹配的时候的游戏段位或者分数，玩家实力（这个也很难说是什么，毕竟胜率，游戏分数都能体现一部分实力），如果实力有波动那胜率还能计算吗？\\
3.在上面情况下，有没有可能造成会有elo的感觉？（什么是elo？多次连败？连胜后连败？）\\
4.脱离elo来讲这和训练模型或许有一定关系，target究竟是胜率还是实力？训练到模型稳定后有没有胜负数据？或许可以直接来说明胜负情况。
\end{quote}

提问者随后追问：厂商常说“每局双方约五五开”，这指的是\emph{实力}五五开，还是\emph{胜率}五五开？此外，提问者要求用实力波动模型展示“上头”，并在合适处说明结论对《炉石传说》酒馆战棋这类 8 人名次游戏意味着什么。

\section*{四项研究与章节的对应}

\begin{itemize}
  \item \textbf{研究一（\cref{ch:lit}）}：文献与一手材料调研。前半部分综述评分模型、匹配机制和各游戏的学术研究；后半部分逐家阅读厂商的开发者文章和白皮书、24 件美国授权专利的权利要求、两起诉讼的案卷与报道，并按证据强度分级。
  \item \textbf{研究二（\cref{ch:winrate}）}：在“只有输赢、匹配不被操纵”的前提下推导胜率。按 1v1 随机匹配、1v1 按分匹配、团队、8 人自由混战和实力波动分节给出闭式公式，并用模拟验证。
  \item \textbf{研究三（\cref{ch:streaks}）}：把玩家口中的“Elo”形式化为可检验的假说，计算纯随机和诚实 Elo 下的连败与“连胜后连败”，用五种实力波动模型（含“上头”）看什么会放大它们，再给出操纵会留下的统计指纹。
  \item \textbf{研究四（\cref{ch:ml}）}：从训练模型的角度重述评分系统：标签、输出与潜参数各是什么，收敛后的胜负数据还有多少信息，怎样用胜负数据做检验，以及自博弈和大模型评测中的同一结构。
\end{itemize}

最后的“结论与展望”把四项研究串成一条主线，并说明两项后续工作：收集社交平台上的“控胜率”说法并用第三章的模型检验，以及分析提问者本人的对局数据。推导与模拟参数集中在附录（\cref{app}）。

\section*{问题速答}
\phantomsection\addcontentsline{toc}{section}{问题速答}\label{sec:intro-qa}

下表把提问者的每个子问题对应到一句话回答和正文位置。带“延伸”字样的是正文顺带回答、提问者未直接问到的问题。

{\footnotesize
\setlength{\tabcolsep}{3pt}
\renewcommand{\arraystretch}{1.18}
\begin{longtable}{@{}>{\raggedright\arraybackslash}p{3.5cm}>{\raggedright\arraybackslash}p{9.6cm}>{\raggedright\arraybackslash}p{2.0cm}@{}}
\caption*{\small\textbf{问题速答}：子问题、一句话回答与正文位置}\\
\toprule
问题 & 一句话回答 & 见 \\
\midrule
\endfirsthead
\toprule
问题 & 一句话回答（续） & 见 \\
\midrule
\endhead
\bottomrule
\endfoot
\multicolumn{3}{@{}l}{\textbf{研究一：文献、披露与专利}}\\
通用的 Elo 或匹配机制有人研究吗？ & 有，且成熟：Elo、Glicko、TrueSkill 等评分模型的统计性质已有严格结果；匹配研究覆盖公平与信息量、等待时间、分队和以留存为目标的匹配。 & \cref{sec:1a-models,sec:1a-mm} \\
各个游戏专门的呢？ & 有但不均衡：只有微软系游戏公开了完整评分模型，其余多为外部研究者在公开数据上的行为研究；酒馆战棋没有评分或匹配研究。 & \cref{sec:1a-games} \\
游戏公司有披露吗？ & 有：普遍是“隐藏分负责匹配、可见段位另算”，承认一组非实力规则（延迟、新手与回流、组排、补位、行为分、机器人），否认按连胜连败安排队友。 & \cref{sec:ch1b-disclosure} \\
有公司的专利或论文吗？ & 有：EA 的留存匹配专利家族、Activision 的“诱导购买”匹配专利，以及 EOMM、EnMatch、CUPID 等论文；专利不证明部署，论文的部署游戏不具名。 & \cref{sec:ch1b-patents,sec:1a-mm-engage} \\
\midrule
\multicolumn{3}{@{}l}{\textbf{研究二：不干涉匹配时的胜率}}\\
只有输赢、不干涉匹配，能算出胜率吗？ & 能算期望：随机匹配下 $W=\Phi\bigl((s-\mu)/\sqrt{\sigma_d^2+\tau^2}\bigr)$；按分匹配下由恒等式“胜率 $=$ 平均预测胜率 $+$ 评分净变化$/(K\times$局数$)$”决定，长期趋于 50\%。 & \cref{sec:ch2-random,sec:ch2-identity} \\
和一局的人数有关吗？ & 有：团队中个人优势约按 $1/\sqrt n$（加和模型）到 $1/n$（均值模型）稀释；8 人局的基准率取决于“胜”的定义（第一名 12.5\%，前四 50\%）。 & \cref{sec:ch2-team,sec:ch2-ffa} \\
和匹配时的段位或分数有关吗？ & 只通过对手分布起作用：评分落后于实力时胜率偏高（重置为 1500 的高手前 50 局约 77.5\%），天梯两端因人少偏离 50\%；可见段位叠加了保底和加分，不宜直接用来估计。 & \cref{sec:ch2-window,sec:ch2-reset,sec:ch2-rank} \\
“实力”是什么？胜率、分数各体现多少？ & 实力是只在比较中有意义的潜变量，可操作地定义为“对固定参照池的胜率”；胜率是实力相对实际对手的位置，评分是带噪声、滞后和刻度的估计。 & \cref{sec:ch2-defs,sec:ch4-identify} \\
实力有波动，胜率还能算吗？ & 能算期望，但误差更大：独立波动让曲线变平；自相关的状态让百局胜率的方差变为二项分布的约 1.6–3.6 倍（模拟）；持续进步让胜率高出预测约 $\lambda v/K$。 & \cref{sec:ch2-fluct} \\
（延伸）多少局才能看出差别？ & 区分 52\% 与 50\% 约需 4,900 局（双侧 5\%、检验力 80\%）；8 人局用完整名次比用前四率省样本。 & \cref{sec:ch2-stats} \\
\midrule
\multicolumn{3}{@{}l}{\textbf{追问：每局“五五开”}}\\
每局双方约五五开，是实力五五开还是胜率五五开？ & 是系统\emph{预测}的胜率五五开，即双方\emph{估计}实力相等；真实实力并不相等（$K=16$ 时 90\% 的“五五开”对局真实胜率在 38\%–62\% 之间），长期胜率接近 50\% 是收敛的结果，多家厂商明说不是目标。 & \cref{sec:ch2-fair,sec:ch1b-fifty} \\
\midrule
\multicolumn{3}{@{}l}{\textbf{研究三：“Elo 感”}}\\
上述情况下会有“Elo 的感觉”吗？ & 会，不需要操纵：胜率五成的玩家打 100 局，81\% 会遇到 5 连败；诚实的按分匹配带来的均值回归单局很小，按 20–50 局的战绩看约 $\pm4$–5 个百分点（1v1）。 & \cref{sec:ch3-random,sec:ch3-elo} \\
什么是 Elo？ & 学术上是一套评分更新公式；玩家说的“ELO 机制”指被怀疑的操纵（强制五成、连胜惩罚、Elo 地狱），本文把它们写成可检验的假说。 & \cref{sec:1a-elo,sec:ch3-hyp} \\
多次连败？ & 很常见：100 局中最长连败的期望约 6 局，1000 局内出现 10 连败的概率约 39\%；诚实 Elo 下连败使对手变弱，长连败反而略少；团队游戏的“出不去”是评分追得慢，不是卡住。 & \cref{sec:ch3-longest,sec:ch3-elo} \\
连胜后连败？ & 纯随机下 100 局内出现“4 连胜紧接 4 连败”的概率为 31\%；诚实 Elo 只让 5 连胜后下一局胜率低约 3 个百分点（1v1，$K=32$）；真正放大它的是连败触发的“上头”（“5 连胜紧接 5 连败”约翻倍）。 & \cref{sec:ch3-wklk,sec:ch3-tilt} \\
“上头”等实力波动怎样建模？ & 写成 $s_t=s+u_t$ 的五种形式（独立噪声、AR(1)、两状态、上头、疲劳）：慢变化大多被评分吸收，只有快速、由失败触发的上头显著增加长连败。 & \cref{sec:ch3-fluct} \\
（延伸）若有操纵，会留下什么痕迹？ & 在匹配输入（连胜后对手评分显著偏高、在触发阈值处跳变，或分队构成）和系统自己的预测里，而不在胜负序列里；对称的“控五成”反而让长连败更少，“输家队列”只有与上头叠加才明显放大连串。 & \cref{sec:ch3-manip} \\
\midrule
\multicolumn{3}{@{}l}{\textbf{研究四：训练视角}}\\
这和训练模型有关吗？ & 有：Elo 更新恰是对数损失上的一步随机梯度下降，$K$ 是学习率；TrueSkill 是学习率随不确定度变化的贝叶斯版本。 & \cref{sec:ch4-sgd} \\
target 是胜率还是实力？ & 评分系统里，标签是胜负，输出是胜率，被学到的潜参数是实力；只有在自博弈中，胜率本身才是优化目标。 & \cref{sec:ch4-roles} \\
训练稳定后还有胜负数据吗？ & 有，照样产生、照样推动评分；只是对负责匹配的模型而言每局接近抛硬币，而五五开的对局恰恰 Fisher 信息最大。 & \cref{sec:ch4-coinflip} \\
能直接用胜负说明胜负情况吗？ & 能，但要把冻结的、校准过刻度的收敛模型当零假设，检验“实际减预测”的残差；朴素的“连败后胜率”检验在诚实 Elo 下也会误报。 & \cref{sec:ch4-null} \\
（延伸）对酒馆战棋 AI 意味着什么？ & 用名次似然训练，用冻结参照池评估；前四率提高 2 个百分点约需 4,900 局，平均名次改进 0.1 约需 4,100 局。 & \cref{sec:ch4-bg} \\
\end{longtable}
}

\section*{记号}

全文统一使用下列记号；只在某一节使用的局部记号在首次出现处定义。

\begin{table}[htbp]
  \centering\small
  \caption*{\small\textbf{记号表}}
  \begin{tabular}{@{}ll@{}}
    \toprule
    记号 & 含义 \\
    \midrule
    $s_i$ & 玩家 $i$ 的真实实力（潜变量，Elo 刻度） \\
    $\mu,\tau$ & 群体实力的均值、标准差（默认 $\tau=200$） \\
    $R_i$ 或 $R_{i,t}$ & 系统对实力的估计（评分、隐藏分、MMR） \\
    $\Elogit(x)=1/(1+10^{-x/400})$ & Elo 胜率曲线 \\
    $\cElo=\ln 10/400$ & Elo 曲线在 0 处的斜率常数，$\Elogit'(0)=\cElo/4$ \\
    $E$ & 系统预测胜率 $\Elogit(R_i-R_j)$ \\
    $p$ & 真实单局胜率 $\Elogit(s_i-s_j)$ \\
    $y\in\{0,1\}$ & 单局结果（Elo 文献中的 $S$） \\
    $\Kf$；$T$ & K 因子；局数 \\
    $n$；$N$ & 每队人数；每局人数（自由混战） \\
    $\sigma_q$ & 排队/匹配噪声（匹配窗口宽度） \\
    $\sigma_e$ & 评分围绕实力的稳态误差标准差 \\
    $\lambda$ & 评分刻度相对实力刻度的伸缩系数 \\
    $u_t$ & 实力波动（状态），$s_{i,t}=s_i+u_t$ \\
    $\Phi$ & 标准正态分布函数 \\
    \bottomrule
  \end{tabular}
\end{table}

\paragraph{局部记号。}下列符号在某一节另有局部含义，首次出现处都有定义：AR(1) 状态的自回归系数在第二、三章都记为 $\rho$，每局评分吸收比例（学习速率）都记为 $b$；第二章的 $r$ 是随机匹配下胜率曲线的压缩系数，$\kappa_R$ 是误差衰减的修正系数；第三章的 $\kappa$ 是“评分误差传到真实胜率的比例”，$\psi$ 是疲劳斜率，$\nu$、$r_0$ 只用于最长连败的 Gumbel 推导；第四章的 $\gamma_s$ 是评分噪声的收缩系数，$\hat\gamma$ 是拟合的刻度；$\Lambda(x)=1/(1+e^{-x})$ 是逻辑斯蒂函数，与表示标准差的 $\sigma$ 区分。TrueSkill 的不确定度在第四章写作 $\varsigma$、在第一章沿用原论文的 $\sigma$，其中的常数 $c$ 写作 $\tilde c$；Glicko-2 的系统常数写作 $\tau_{\mathrm G}$。仍保留的一处重复是 EOMM 一节的 $c(\cdot)$（流失概率），正文已注明。

\paragraph{术语对照。}厂商和社区对同一概念用词不一，本文统一如下（\cref{tab:intro-terms}）。

\begin{table}[htbp]
  \centering\small
  \caption{术语对照}
  \label{tab:intro-terms}
  \begin{tabularx}{\linewidth}{@{}>{\raggedright\arraybackslash}p{3.0cm}>{\raggedright\arraybackslash}X@{}}
    \toprule
    本文用语 & 原文或其他常见说法 \\
    \midrule
    评分 $R$ & MMR、隐藏分、内部评分、匹配分、技能分、原始技能、竞技分、等级分（Elo rating）\\
    可见段位 & 段位、星、段位分、外显评分、RR、SR、CSR \\
    实力 $s$ & 厂商文章中的 skill 有时指实力、有时指评分，按上下文译为“实力 $s$”或“评分 $R$” \\
    组排 & 车队、预组队、开黑（party、premade、squad）\\
    输家队列 & loser queue（玩家所说的“系统让连败者继续输”的匹配）\\
    按分匹配 & 技能匹配（skill-based matchmaking, SBMM）\\
    检验力 & 功效、统计功效（power）\\
    评分差 & “X 分”均指 Elo 刻度上的分差，即 $\Elogit$ 的自变量 \\
    \bottomrule
  \end{tabularx}
\end{table}

\section*{证据标签}

每条事实性陈述后标注来源性质：
\begin{itemize}
  \item \tagLit：同行评审论文或公司技术报告；只有 arXiv 预印本、尚无同行评审版本的，首次出现时注明“预印本”；
  \item \tagOff：开发商官方博客、白皮书、员工以官方身份的公开发言，以及法院案卷、判决、法规文本本身；案卷中\emph{当事人的指控}标\tagCom；
  \item \tagPat：专利文献，只说明申请并获得过保护，不说明已部署；
  \item \tagCom：玩家、自媒体、第三方工具、百科，未经证实；
  \item \tagOwn：本项目自己的推导或模拟，正文或图注给出参数。
\end{itemize}
“无需操纵即可解释”不等于“已证明没有操纵”；全文的措辞都保持这个区分。

\section*{方法与局限}

\paragraph{来源。}检索与阅读集中在 2026 年 10 月 18—19 日。核心来源读的是全文或原页面，包括 TrueSkill 与 TrueSkill 2、Glickman 的评分综述、Aldous 等 Elo 理论论文、EOMM、EnMatch、CUPID、QuickSkill、Kang 等的流失研究、Rosenthal、Sapienza、Chowdhary、Miller–Sanjurjo、AlphaGo Zero 与 AlphaStar 的作者稿、OpenAI Five、Chatbot Arena；Riot、Valve、Blizzard、Activision、Respawn、343、KRAFTON 的官方文章与白皮书；24 件美国授权专利的首页与独立权利要求（USPTO 全文图像，经 OCR）；Zajonc 诉 EA 的联邦法院案卷。只读到摘要、题录或二手转述的来源，正文在对应位置加了“待核对”脚注，并逐条列入仓库的 \texttt{docs/TO\_VERIFY.md}（约 40 条），主要是：1960 年以前的经典论文（只核对了题录）、付费墙后的论文（如 Kou 等、Deng 等、Chen–Elmachtoub–Lei、几篇心理学论文）、知网上的中文学位论文、腾讯 2018 年官网原文（依据媒体摘录）、王者荣耀案判决书（未公开，依据多家媒体报道），以及中国专利的检索（中国知网、万方、百度学术和国家知识产权局的检索页在本研究的网络环境中无法访问：知网证书与主机不符，百度学术返回 403，万方检索页需脚本渲染）。

\paragraph{模型。}本文所有\tagOwn{}数字都来自理想化模型：一维实力，逻辑斯蒂胜率曲线，正态分布的群体（$\tau=200$），固定 $\Kf$ 的 Elo，按“评分加噪声排序后相邻配对”的简化匹配器；没有组排、位置、延迟、版本更新和玩家流失。“上头”“疲劳”等波动模型的参数是示意值，不是从真实数据估计的。因此具体数字会随参数变化，可靠的是公式给出的依赖关系与各机制之间的相对大小。代码与缓存在仓库的 \texttt{code/} 目录，每张图的图注写明了参数。
